\section{Method}
\label{sec:method}

\subsection{Overview}

Our method takes as input a reference set
$\mathcal{R} = \{ (V_i, N_i) \}_{i=1}^{N}$, where $V_i = (V_i^1, \dots, V_i^T)$ is a tactile normal
video of $T$ frames and $N_i$ is the surface normal map rendered at the pose of touch $i$, together
with a normal map $N_q$ rendered at a query sensor pose, which serves as the geometric context for
prediction. Our goal is to produce the predicted video $\hat{V}_q$. Obtaining these renderings does
not require an object mesh at deployment time: one can differentiate a depth map from a commodity
depth or stereo camera into camera-frame normals and project them into the sensor frame using the
sensor-to-camera pose, known from forward kinematics or a fiducial marker on the sensor. We take a coarse-to-fine approach with three steps: retrieving
the most promising reference touch with DINOv3 features (Sec.~\ref{sec:retrieval}), coarsely
aligning the two normal maps to obtain a homography that gives an initial transfer of the reference
video (Sec.~\ref{sec:coarse}), and refining that transfer with a neural network
(Sec.~\ref{sec:refine}).

\subsection{Retrieval}
\label{sec:retrieval}

A reference touch is useful only if the surface it was taken on resembles the surface at the query
pose, and we measure this resemblance in the space of a self-supervised visual descriptor. We render
normal maps over a field of view four times the sensor footprint, so that the descriptor sees the
surrounding context and not just the small contact patch, and extract the $\ell_2$-normalized
class-token feature $f=\phi(N)/\lVert \phi(N) \rVert_2$ of a pre-trained DINOv3 encoder
$\phi$~\cite{dinov3}. The retrieved reference is the one most similar to the query in the cosine
sense,
\begin{equation}
  i^{*} \;=\; \arg\max_{i \in \{1,\dots,N\}} \; f_i^{\top} f_q .
  \label{eq:retrieval}
\end{equation}
We use the top-1 result throughout. As the descriptor is frozen and the reference features can be
cached, retrieval reduces to a single matrix--vector product at query time.

\subsection{Coarse Alignment}
\label{sec:coarse}

Retrieval tells us which reference to use, and alignment tells us how to place it. On
the top-1 retrieval $N_{i^{*}}$ we run a local feature matcher, SuperPoint
keypoints~\cite{superpoint} matched by SuperGlue~\cite{superglue}, against $N_q$, obtaining $K$
pixel correspondences $\{(\mathbf{p}_k, \mathbf{q}_k)\}_{k=1}^{K}$ from which we estimate a
homography $\mathbf{H} \in \mathbb{R}^{3\times 3}$ under RANSAC~\cite{ransac},
\begin{equation}
  \mathbf{H} \;=\; \arg\min_{\mathbf{H}} \sum_{k=1}^{K}
  \rho\!\left( \bigl\lVert \pi\bigl(\mathbf{H}\,\tilde{\mathbf{p}}_k\bigr) - \mathbf{q}_k
  \bigr\rVert_2 \right),
  \label{eq:homography}
\end{equation}
where $\tilde{\mathbf{p}}_k$ is $\mathbf{p}_k$ in homogeneous coordinates, $\pi$ is the perspective
division, and $\rho$ is the RANSAC inlier--outlier cost. Applying $\mathbf{H}$ to every frame of
the retrieved reference video, $W^{t} = \mathcal{W}( V_{i^{*}}^{t}; \mathbf{H} )$ with
$\mathcal{W}$ denoting image warping, gives the initial transfer. A single transform is shared by
all frames, since the sensor pose does not move relative to the surface during a press, so what
changes across the video is the gel deformation, not the alignment.

\subsection{Learning-based Refinement}
\label{sec:refine}

The initial transfer is geometrically plausible but may be physically inaccurate: it carries the gel
dynamics of the reference press, inherits any residual misalignment left by
Eq.~\ref{eq:homography}, and knows nothing about the query surface geometry beyond what survived
the alignment. Our final step corrects these artifacts with a neural network, balancing accuracy
against the efficiency AR/VR and robotics demand.

\paragraph{Architecture}
Our refinement network is based on ReBotNet~\cite{rebotnet}, a lightweight video enhancement
architecture that processes a pair of consecutive frames through a ConvNeXt-style~\cite{convnext}
spatio-temporal branch and a per-frame token branch, merges them in a shared decoder, and adds the
current frame back through a global residual connection. As it was designed to remove blur and
compression artifacts from natural video, we adapt it in two ways.

\emph{Normal map concatenation.} We concatenate the query normal render $N_q$ to the input as three
aligned channels, broadcast across both frames of the pair. This makes the network aware of the
surface geometry at the query location, which the warped reference cannot supply.

\emph{Temporal conditioning with FiLM.} The normal map is rendered once at the query pose and does
not change over the press, so on its own it tells the network nothing about where in the press it
currently is. We therefore inject the frame's normalized timestamp $\tau_t = (t-1)/(T-1) \in
[0,1]$, which runs from no contact through deepest press to take-off. Following the standard
diffusion-model recipe~\cite{ldm}, a sinusoidal embedding of $\tau_t$ is passed through a small
multi-layer perceptron and mapped to a per-stage scale and shift, applied as feature-wise linear
modulation (FiLM)~\cite{film}:
\begin{equation}
  \mathrm{FiLM}(\mathbf{h}_s; \tau_t) \;=\;
  \bigl(1 + \gamma_s(\tau_t)\bigr) \odot \mathbf{h}_s + \beta_s(\tau_t),
\end{equation}
where $\mathbf{h}_s$ are the features at stage $s$. The heads producing $\gamma_s$ and $\beta_s$
are zero-initialized, so the module starts as an exact identity and departs from it only as
training justifies. Writing the refinement network as $\mathcal{F}$, the prediction for frame $t$
is $\hat{V}_q^{\,t} = \mathcal{F}( [ W^{t-1}, W^{t} ] \oplus N_q ; \tau_t )$, with $\oplus$
denoting channel concatenation and $W^{0} := W^{1}$ for the first frame.

\paragraph{Training loss}
We supervise with a Charbonnier loss $\mathcal{L}_{\mathrm{ch}} = \frac{1}{|\Omega|}\sum_{x \in
\Omega} \sqrt{(\hat{V}_q^{\,t}(x) - V_q^{t}(x))^2 + \epsilon^2}$, a smooth approximation to the
$\ell_1$ distance that is less sensitive to outliers than a squared error. It is evaluated only on
the region $\Omega = \{ x : \max_{c} \lvert V_q^{t}(x, c) - W^{t}(x, c) \rvert > \delta \}$ where
the ground-truth frame differs from the warped input, with the maximum running over the three
normal channels and $\delta = 10^{-2}$. Pixels at which the warped input already agrees with the
ground truth carry no training signal, and keeping them would let the large, unchanged background
dominate the loss.
